\section{An exact refinement of the cyclotomic annihilator}
\label{sec:cyc}
An intermediate step of the readable edition constructs an integer $X$ from a
Bernoulli matrix and proves that $4X^2$ kills a family of Galois cocycles.
Here we solve an optimization problem inside this argument. The value of a
single integer meeting the Lean hypotheses and the effective integer obtained
by combining several such proofs can differ. At modulus five and weight one
we determine both exactly. The statements in this section use the cyclotomic
lemmas imported by the proof. They are independent of zero location in the
critical strip.

\subsection{The upstream interface}
Fix integers $N,e\ge1$. Put $U=(\mathbb Z/N\mathbb Z)^\times$ and use the
representatives $0\le t<N$. Write
\[
 b(t)=(e+2)!N^{e+1} B_{e+1}(t/N),\quad K=(e+2)!(e+1)N,\qquad
 F(c)=(c^e-1)(c^{e+1}-1).
\]
Here $B_j(x)$ is the Bernoulli polynomial with $B_1(x)=x-1/2$.
The values $b(t)$ are integers. Call a character relevant if
$\chi(-1)=(-1)^{e+1}$, and put
$B_\chi=\sum_{t\in U} b(t)\chi(t)$. These sums are nonzero for all relevant
characters by the functional equation and nonvanishing at $1+e$.
The principal character for $N=1$ uses $t=0$. Characters at other moduli
extend by zero on nonunits.

The source's predicate \texttt{IsStickelbergerBound} asks a positive
candidate $X$ to satisfy two requirements.
First, for every rational prime $p$, some $c\equiv1\pmod N$ coprime to $p$
has $F(c)\mid X$. Such a $c$ must be at least two for positive $X$.
Second, for each $c\equiv1\pmod N$ and each relevant character
in a characteristic-zero integral domain $O$, there must be $w\in O$
independent of the modulus $m$ such that
\begin{equation}
 (c^{e+1}-1)X\equiv 2w S_{m,c,\chi}\pmod{mO},\quad
 S_{m,c,\chi}=\sum_{0\le a<Nm}\left\lfloor\frac{ca}{Nm}\right\rfloor
 a^e\chi(a),\quad (m,c)=1.                         \label{eq:stick}
\end{equation}
The source proves
$Kc^eS_{m,c,\chi}\equiv(c^{e+1}-1)B_\chi\pmod{mO}$.
It chooses the matrix
\[
 T_{u,v}=2b(t)+(-1)^e\mathbf1_{t=1}+\mathbf1_{t=-1},\qquad
 t=vu^{-1}\pmod N ,
\]
and supplies the witness
$X_0=F(1+N)F(1+2N)|\det T|$. Character vectors diagonalize $T$.
Relevant characters have eigenvalue $2B_\chi$, and the others have
eigenvalue $2(-1)^e$. The adjugate clears every relevant eigenvalue at once.
See \src{Lib/Cyclotomic/Voronoi.lean}{Voronoi, lines 134--151, 245--256 and 384--459}.

\subsection{The common denominator of the character values}
For each complex relevant $\chi$, let $d_\chi$ be the order of its image
and let $D_\chi$ generate the positive ideal
$(2B_\chi)\mathbb Z[\zeta_{d_\chi}]\cap\mathbb Z$.
Set $D=\mathop{\rm lcm}_\chi D_\chi$, with empty lcm one.
Equivalently $D_\chi$ is the least common denominator of the rational
coefficients of $(2B_\chi)^{-1}$ in the cyclotomic power basis. This gives a
finite exact calculation for every $N,e$.

\begin{proposition}
$D$ is the least positive integer divisible by every relevant $2B_\chi$ in
every characteristic-zero character value domain. For any coprime $a,b\ge2$
congruent to $1\pmod N$,
\[
             X(a,b)=\mathop{\rm lcm}(D,F(a),F(b))
\]
satisfies both upstream requirements.
\end{proposition}
\begin{proof}
The finite image of any such character generates a cyclotomic subring of
$O$. Indeed, in its fraction field the image is cyclic, and a primitive
$d$th root in characteristic zero has minimal polynomial $\Phi_d$ over
$\mathbb Q$. Thus each complex-character divisibility transfers to $O$.
Taking $O=\mathbb Z[\zeta_{d_\chi}]$ proves the converse. An algebraic
integer in $\mathbb Q(\zeta_d)$ belongs to $\mathbb Z[\zeta_d]$, so one
may equally compute the denominators in a single cyclotomic field containing
all character values.
For $X=X(a,b)$ the quotient $u=X/(2B_\chi)$ lies in $O$.
Multiply the upstream congruence by $2u$ to obtain \eqref{eq:stick} with
$w=uKc^e$. Each rational prime is coprime to at least one of $a,b$,
giving the first requirement. The adjugate argument also shows $D\mid|\det T|$.
\end{proof}

\subsection{An exact local calculation}
Define $M=\prod_p p^{m_p}$ using the following finite prescription.
Write $v_p$ for the ordinary integer valuation. For odd $p$,
\begin{equation}
m_p=\begin{cases}
2v_p(N)+v_p(e(e+1)) & p\mid N,\\
\displaystyle\sum_{j\in\{e,e+1\},\ p-1\mid j}(1+v_p(j)) & p\nmid N.
\end{cases}                                                    \label{eq:mp}
\end{equation}
For $p=2$ put $m_2=v_2(e(e+1))+2v_2(N)$ if $4\mid N$, and
$m_2=v_2(e(e+1))+3$ otherwise. Set $H=\mathop{\rm lcm}(D,M)$.

\begin{theorem}\label{th:gcd}
The gcd of the integers $X(a,b)$ over all admissible coprime pairs is exactly
$H$. Every positive integer satisfying the upstream first requirement is a
multiple of $M$.
\end{theorem}
\begin{proof}
For each $p$ minimize $v_p(F(c))$ over $c\ge2$ with $c\equiv1\pmod N$
and $p\nmid c$. For odd $p\mid N$, the lifting-the-exponent formula gives
$2v_p(c-1)+v_p(e(e+1))$ and one can take $v_p(c-1)=v_p(N)$.
For odd $p\nmid N$, an exponent divisible by $p-1$ forces at least
$1+v_p(j)$ in $c^j-1$ by Fermat's theorem and lifting the exponent.
A primitive root modulo $p^2$, combined with the progression modulo $N$,
attains equality for all $j$ at once and contributes zero for the other
exponents. At two, let $a'=v_2(c-1), b'=v_2(c+1)$. The two exponents have
opposite parity, giving total valuation
$2a'+b'-1+v_2(e(e+1))$. If $4\mid N$ we can take
$a'=v_2(N), b'=1$. Otherwise $c\equiv3\pmod8$ is possible and minimizes
$2a'+b'$ at four. This proves the formulas and $M\mid X$ for every
positive $X$ passing the first requirement.

Fix $p$ and choose one minimizing $a$. Choose $b>1$ congruent to $a$ modulo
$p^{m_p+1}$ and to $1$ modulo $N$ and every prime divisor of $a$. These
congruences are compatible, including when $p\mid N$, and give
$(a,b)=1$ and $v_p(F(b))=m_p$. Thus
$v_p(X(a,b))=\max(v_p(D),m_p)$ is attained for each $p$.
Taking the gcd proves the claim. A finite set of witnesses always suffices:
start with any one witness and minimize at each of its prime divisors.
\end{proof}

\begin{corollary}\label{cor:kill}
In the setting of the repository's odd component with coefficients
$O/\ell^\nu$, the integer $4H^2$ annihilates all its cocycles. For
$\ell>N+e+3$, at most $v_\ell(H)+1$ Galois elements suffice to detect an
element by evaluation. For a CM subfield $E\subseteq\mathbb Q(\mu_N)$,
under the source's coefficient hypotheses, $16\varphi(N)H^2$ annihilates
the module \texttt{shaOdd} at weight $e$ and every level.
\end{corollary}
\begin{proof}
For each genuine witness $X$, the source proves the assertions with $X$ in
place of $H$. Take the gcd of the integer annihilators using B\'ezout's
identity. For detection choose a witness attaining the minimum valuation at
the specified $\ell$. These two operations are different. In particular
the proof never requires $H$ itself to pass the first requirement.
See \src{Lib/Cyclotomic/Kolyvagin/ReciprocityComponentBounds.lean}{component
annihilation and detection, lines 676 and 1104}, and
\src{Lib/Cyclotomic/OddSha.lean}{CM descent, line 535}.
\end{proof}
Here $O$ is the integer ring of a finite extension of $\mathbb Q_\ell$ as in
the source. The component consists of locally constant cocycles on
$G_\mathbb Q$ for $(\mathrm{cyc}^e\chi)^{-1}$, of odd parity, subject to its
specified vanishing at conjugation, inertia and the local set at $\ell$.
CM descent also asks for enough roots of unity in $O$. The finite-level
CM object uses cocycles with its particular local conditions.
The identification of this object with a conventional cohomology group,
and control of even class-group eigenspaces in Vandiver's conjecture,
would require further arguments. The upstream
\src{Lib/Cyclotomic/OddSha.lean}{module overview} states the comparison
issue explicitly. Classical higher Stickelberger theory already gives arithmetic annihilation results (see \href{https://archive.numdam.org/articles/10.1016/j.crma.2003.09.019/}{Banaszak, \emph{Higher analogues of Stickelberger's theorem}}). Our sharpness statement concerns the
displayed construction.

\subsection{The smallest witness at modulus five}
At $N=5,e=1$, the four $b(t)$, in increasing order of $t$, are
$1,-11,-11,1$. The trivial and quadratic characters give relevant
eigenvalues $-40$ and $48$. The two odd characters give $-2$ each.
Thus $D=240$, $|\det T|=7680$, $F(6)=175$, $F(11)=1200$ and
$X_0=1\,612\,800\,000$.

\begin{theorem}
At $(N,e)=(5,1)$ the least positive integer satisfying the full upstream
predicate is $8400$. The gcd of \emph{all} its positive witnesses is $1200$.
\end{theorem}
\begin{proof}
The pair $(6,11)$ gives $8400$ and $(11,16)$ gives $61200$ since
$F(16)=3825$. Their gcd is $1200=M$, a divisor of every witness by
Theorem~\ref{th:gcd}. For $0<X\le8400$, the only possible auxiliary
integers are $6,11,16$: $F$ is increasing and $F(21)=8800$.
The eligible integers with $F(c)\mid X$ must have gcd one. The only way
within the budget is to include $6$ and $11$: $6$ and $16$ share a
factor, and $\mathop{\rm lcm}(F(11),F(16))=61200$.
Thus $8400\mid X$.
\end{proof}
The example proves more than a reduction in the written size of a bound.
The original witness permits the exceptional prime seven. Since
$7\nmid4\cdot1200^2$, the corresponding odd components vanish at seven at
\emph{every} finite level. More generally at modulus five and weight one
only $2,3,5$ can support these components. This is an exclusion of primes;
it does not assert a nonzero component at any of those three.

\begin{center}
\begin{tabular}{rr r r r r}
\toprule
$N$ & $e$ & $D$ & $M$ & $X(1+N,1+2N)$ & $H$\\
\midrule
5&1&240&1200&8400&1200\\
5&2&1440&1200&41\,176\,800&7200\\
7&2&16\,128&2352&1\,986\,179\,328&112\,896\\
8&1&192&384&23\,040&384\\
12&2&17\,280&864&2\,973\,559\,680&17\,280\\
15&1&2880&3600&489\,600&14\,400\\
\bottomrule
\end{tabular}
\end{center}
The checker in \path{checks/cyclotomic.py} uses polynomial inversion modulo
cyclotomic polynomials over $\mathbb Q$. It exhibits finite lists of
coprime-pair witnesses giving each claimed gcd. In the classical scalar
case $(N,e)=(1,11)$ it gives
$H=2^{10}3^4\cdot5\cdot7\cdot11\cdot13\cdot691$.
Consequently at primes greater than $15$ only $691$ survives this test.
It comes from the numerator of $B_{12}=-691/2730$.
This recovers the familiar Bernoulli obstruction in this normalization.
